% Compile from this directory with: pdflatex project_7.tex (twice).
\documentclass[11pt,a4paper]{article}
\usepackage[T1]{fontenc}
\usepackage[utf8]{inputenc}
\usepackage[margin=22mm]{geometry}
\usepackage{lmodern}
\usepackage{amsmath,amssymb}
\usepackage{enumitem}
\usepackage{xcolor}
\usepackage{microtype}
\usepackage{hyperref}
\definecolor{epflred}{HTML}{E3000B}
\hypersetup{colorlinks=true,urlcolor=epflred,linkcolor=epflred,
  pdftitle={ENG-654 Project 7: Signed IK Classification of PUMA 560 and UR5}}
\setlength{\parindent}{0pt}
\setlength{\parskip}{6pt}
\setlist[itemize]{leftmargin=*,itemsep=5pt,topsep=4pt}
\urlstyle{tt}

\begin{document}
{\small\textbf{EPFL \textbar\ ENG-654}\hfill Kinematics-Grounded Motion Planning for Robots}
\vspace{5mm}

{\color{epflred}\Large\bfseries Project 7}

{\LARGE\bfseries Signed IK Classification of PUMA 560 and UR5\par}

\textbf{Format:} Group of two students; simulation-based project.\\
\textbf{Course foundations:} Lectures 2, 4--6; Exercises 2--3.

\section*{Motivation}
A useful IK label describes a geometric configuration and the singular
boundary separating it from alternatives. PUMA 560 and UR5 offer two different
6R architectures for testing this idea. Comparing them shows how determinant
factorization can support an interpretable classification without assuming
that identical joint numbers imply identical geometry.

\section*{Task}
\textbf{Overall objectives.} Compare the PUMA and UR5 architectures, derive and interpret the three
factors of the UR5 Jacobian determinant, associate regular IK solutions with
signed categories, and deduce whether UR5 is cuspidal or noncuspidal.

\begin{itemize}
  \item \textbf{Compare the two kinematic models.} Inspect the URDFs, STL geometry,
axis arrangements, and D--H conventions. Identify the PUMA spherical wrist and
the different UR5 wrist/parallel-axis structure. Validate both forward models
and document joint offsets, base frames, and tool frames.

  \item \textbf{Derive and check the UR5 determinant.} Construct the full geometric
Jacobian and factor its determinant into three configuration-dependent
components, ignoring nonzero constant factors. Verify the factorization
numerically against direct Jacobian evaluation and interpret each zero set
geometrically.

  \item \textbf{Assign signed categories.} Use the provided IK routines to obtain
eight regular mathematical solutions at a suitable target for each robot.
Assign $(\operatorname{sgn}f_1,\operatorname{sgn}f_2,\operatorname{sgn}f_3)$
to every UR5 solution and compare with PUMA's shoulder--elbow--wrist labels.
Check the one-to-one association at additional regular targets and separate
joint-limit filtering from mathematical classification.

  \item \textbf{Visualize the separators.} Plot UR5 determinant factors on selected
$(q_2,q_3)$ slices at explicitly fixed $q_4,q_5$ and other joint values.
Relate these slices to full six-dimensional configurations without treating
projected off-slice IKs as on-slice points. For the PUMA arm, also map reduced
singularities and the two-/four-arm-IK workspace distribution.

  \item \textbf{Deduce cuspidality from structure.} Explain why changing a factor
sign along a continuous path requires a zero crossing. Combine this fact with
the complete regular IK classification to infer UR5's cuspidality status;
factor count or determinant sign alone is not a sufficient argument. Test
small approaches to each separator and report ambiguous near-zero labels.
\end{itemize}

\newpage
\section*{Specifics and scope}
Study two architectures, three regular targets per robot, and one short
approach to each UR5 singularity type. The UR5 URDF, STL meshes, and IK
implementation will be provided. Reuse the PUMA IK baseline; focus new derivation
on the UR5 Jacobian and classification rather than implementing another general
IK solver. Allow geometric sign conventions to differ between models.

Check both translation and orientation residuals
for every claimed full-pose IK. Document angle periodicity and limits,
Jacobian conventions, and any scaling used for singular values. State all
fixed coordinates in reduced plots; a two-dimensional slice does not show
the entire 6R singular set. Refine decisive transitions, and distinguish
sampled evidence from a continuous-path guarantee. Collision freedom is
not implied by these kinematic checks.

\section*{Expected outcome}
Explain both robot geometries and show how each UR5 IK fits into a signed
category of three determinant factors. Deliver the verified factorization,
illustrated classification tables, singularity plots, and a justified deduction
of whether UR5 is cuspidal or noncuspidal.

Submit reproducible code, model parameters and test data, the required plots,
a concise 5--6-page report, and a short simulation demonstration. Explain
both successful cases and failures; conclusions must follow the numerical
and geometric evidence rather than an expected outcome.

\section*{Resources}
The PUMA 560 URDF is
\path{lectures_main/assets/models/puma/puma560_robot.urdf}; its STL meshes are
in the same directory. The lecture website provides axis/frame inspection,
IK visualization, and path demonstrations. Confirm joint offsets and the
selected terminal frame when importing the model.
UR5 resources will be supplied separately; use the provided version consistently
instead of substituting another UR-series robot or dimensional variant.


Supporting reading:
\href{https://salunkhedurgesh.com/projects/main/resources/papers/ijrr24_salunkhe.pdf}{\emph{Kinematic issues in 6R cuspidal robots}},
especially the discussion of determinant factors and UR-type architectures.
Explain the argument in your own model conventions rather than copying a
factorization without checking it.

Use the lecture website to verify D--H parameters, visualize IKs and
singularities, and simulate paths. All listed paths are relative to the
repository root. Start the website following \path{lectures_main/README.md}.

\section*{Workload and collaboration}
Budget approximately 60 hours per student (120 person-hours per pair)
for this 2-ECTS project, following EPFL's 30-hours-per-credit convention.\footnote{\href{https://www.epfl.ch/education/teaching/teaching-guide-2/getting-started/design-a-course_1/}{EPFL: Design a Course, workload guidance}.}
Deduct any lectures or exercises included in those credits. Allocate roughly
20 team-hours to model validation, 50 to the core work, 30 to experiments,
and 20 to reporting. Reuse the specified IK and simulations. Share validation
and interpretation even if modelling and implementation are divided between
students. Hardware execution is outside scope.

\end{document}
